% Vietnamese translation for OI Public Library
% Source-File: IOI中国国家候选队论文/2009/高逸涵/数位计数问题解法研究.ppt
\documentclass[11pt]{article}
\usepackage{oipl-vn}

\newcommand{\Oh}{\mathcal{O}}

\title{Nghiên cứu phương pháp giải bài toán đếm theo chữ số\\{\large Ghi chú theo slide}}
\author{Gao Yihan\\Trường Phổ thông trực thuộc Đại học Thanh Hoa}
\date{2009}

\begin{document}
\maketitle

\origtitle{Methods for digit-counting problems}
\sourcepath{IOI-China-Candidate-Team-Papers/2009/Gao-Yihan/digit-counting-slides.ppt}

\section{Chủ đề}

Slide tập trung vào các bài đếm theo chữ số: đếm tổng chữ số, đếm số theo
thứ tự đặc biệt, hoặc tìm thứ hạng trong tập các biểu diễn. Điểm cốt lõi là
tách số thành tiền tố đã cố định và phần hậu tố còn tự do.

\section{Ví dụ}

\begin{itemize}
  \item Tổng chữ số theo vị trí: xét đóng góp của từng vị trí qua chu kỳ của
  chữ số ở vị trí đó.
  \item \textit{The Sum}: dùng đếm theo tiền tố để xử lý tổng trên một đoạn
  số lớn.
  \item \textit{Graduated Lexicographical Ordering}: kết hợp độ dài, trọng số
  và thứ tự từ điển để xác định thứ hạng hoặc phần tử cần tìm.
\end{itemize}

\section{Khuôn mẫu giải}

Với giới hạn \(N\), đọc chữ số từ cao xuống thấp. Ở mỗi vị trí, các lựa chọn
nhỏ hơn chữ số hiện tại tạo ra một khối hậu tố tự do có thể đếm nhanh; lựa
chọn bằng chữ số hiện tại giữ ta trên biên \(N\). Với bài tìm phần tử thứ
\(K\), thủ tục đếm khối được dùng để quyết định chữ số kế tiếp.

\section{Ghi chú}

Bài của Gao Yihan gần với bài của Lưu Thông nhưng nhấn mạnh hơn vào công thức
đóng góp theo vị trí và thứ tự đếm. Đọc hai bài cạnh nhau giúp thấy hai cách
nhìn bổ sung cho cùng họ bài toán.

\end{document}
